\subsection{MAXIMAL TORI OF SUBGROUPS AND QUOTIENT GROUPS}

PROPOSITION 1. Let G and G' be two connected compact Lie groups.

\begin{enumerate}
\def\labelenumi{\alph{enumi})}
\item
  Let \(f: \mathbf{G} \to \mathbf{G}'\) be a surjective morphism of Lie groups. The maximal tori of \(\mathbf{G}'\) are the images under \(f\) of the maximal tori of \(\mathbf{G}\) . If the kernel of \(f\) is central in \(\mathbf{G}\) (for example discrete), the maximal tori of \(\mathbf{G}\) are the inverse images under \(f\) of the maximal tori of \(\mathbf{G}'\) .
\item
  Let \(\mathrm{H}\) be a connected closed subgroup of \(\mathrm{G}\) . Every maximal torus of \(\mathrm{H}\) is the intersection with \(\mathrm{H}\) of a maximal torus of \(\mathrm{G}\) .
\item
  Let H be a connected closed normal subgroup of G. The maximal tori of H are the intersections with H of the maximal tori of G.
\item
  Let T be a maximal torus of G; then L(T) is a Cartan subalgebra of L(G) (no. 2, Th. 2 a)), so L(f(T)) is a Cartan subalgebra of L(G') (Chap. VII, §2, no. 1, Cor. 2 of Prop. 4); it follows that f(T) is a maximal torus of G' (no. 2, Th. 2 a)). If Ker f is central in G, it is contained in T (Cor. 2 of Th. 2), so T = \(f^{-1}(f(T))\) .
\end{enumerate}

Conversely, let \(T'\) be a maximal torus of \(G'\) ; we show that there exists a maximal torus T of G such that \(f(\mathrm{T}) = \mathrm{T}'\) . Let \(T_{1}\) be a maximal torus of G; then \(f(\mathrm{T}_{1})\) is a maximal torus of \(G'\) and there exists \(g' \in G'\) such that \(\mathrm{T}' = g' f(\mathrm{T}_{1}) g'^{-1} (\mathrm{Th.~2~b})\) ; if \(g \in G\) is such that \(f(g) = g'\) , we have \(\mathrm{T}' = f(\mathrm{T})\) with \(\mathrm{T} = g\mathrm{T}_{1}g^{-1}\) .

\begin{enumerate}
\def\labelenumi{\alph{enumi})}
\setcounter{enumi}{1}
\item
  Let S be a maximal torus of H; this is a torus of G so there exists a maximal torus T of G containing S. Then \(T \cap H\) is a commutative subgroup of H containing S, hence is equal to S (no. 2, Remark 2).
\item
  By §1, no. 3, Prop. 2 c), L(G) is the direct product of L(H) with an ideal; thus, the Cartan subalgebras of L(H) are the intersections with L(H) of the Cartan subalgebras of L(G). Thus, for any maximal torus T of G, T∩H contains a maximal torus S of H and S = T∩H (no. 2, Remark 2).
\end{enumerate}

Remarks. 1) Proposition 1 generalizes immediately to connected groups with compact Lie algebras. In particular, if G is a connected Lie group whose Lie algebra is compact, the Cartan subgroups of G (cf.~Remark 3, no. 2) are exactly the inverse images of the maximal tori of the connected compact Lie group Ad(G) (under the canonical homomorphism from G to Ad(G)).

\begin{enumerate}
\def\labelenumi{\arabic{enumi})}
\setcounter{enumi}{1}
\tightlist
\item
  Let G be a connected compact Lie group, \(\widetilde{\mathrm{D}}(\mathrm{G})\) the universal covering of the group \(\mathrm{D}(\mathrm{G})\) and \(f : \widetilde{\mathrm{D}}(\mathrm{G}) \to \mathrm{G}\) the composite of the canonical morphisms from \(\tilde{\mathrm{D}}(\mathrm{G})\) to \(\mathrm{D}(\mathrm{G})\) and from \(\mathrm{D}(\mathrm{G})\) to G. Then the map \(T \mapsto f^{-1}(\mathrm{T})\) is a bijection from the set of maximal tori of G to the set of maximal tori of \(\tilde{\mathrm{D}}(\mathrm{G})\) ; the inverse bijection associates to a maximal torus \(\tilde{T}\) of \(\tilde{\mathrm{D}}(\mathrm{G})\) the maximal torus \(\mathrm{C}(\mathrm{G})_{0}.f(\tilde{\mathrm{T}})\) of G.
\end{enumerate}

